% GENERATED FILE. Edit canonical lessons, then run npm run generate:books.
% canonical-source-hash: fnv1a64:be7172347e1a48a5
% canonical-lessons: GE-C524-komoedie, GE-C524-tragoedie, GE-C524-dokumentarfilm, GE-C524-untertitel, GE-C524-synchronisieren

\chapter{Komödie und Tragödie}
\label{ch:ge-komodie-und-tragodie}
\hlchaptermodality{524}

\begin{chapteropening}
\textbf{By the end of this chapter:} \emph{I can talk about kinds of films.}

Complete the last lesson of chapter 524: I can talk about kinds of films.
\end{chapteropening}

\section[die Komödie]{die Komödie — comedy}

\label{lesson:GE-C524-komoedie}



\begin{quote}
Before the new one: say the German for season, then the German for ratings.
\end{quote}

\subsection*{You'll want to know: die Komödie}
\textbf{die Komödie} — \textquotedblleft{}comedy\textquotedblright{}.

\begin{quote}
Gestern haben wir eine lustige Komödie gesehen.
\end{quote}

\begin{quote}
\emph{Yesterday we saw a funny comedy.}
\end{quote}

Plural: \emph{die Komödien}.

\begin{cousinweb}
\emph{Komödie} comes from Latin \emph{cōmoedia}, from Greek \emph{kōmōidía} \textquotedblleft{}comedy, a play\textquotedblright{} (EtymDB).
\end{cousinweb}

\begin{grammarlens}[title={What you've built}]
One more word for telling, explaining and describing at B1.
\end{grammarlens}

\subsection*{Your turn}
\begin{itemize}
\raggedright
  \item \emph{Say it:} \emph{die Komödie}
  \item \emph{Say it:} \emph{die Komödie}, once more
  \item \emph{Say it:} a sentence of your own with \emph{die Komödie}
  \item \emph{Recall:} say the German for season, then the German for ratings, then say \emph{die Komödie} again
\end{itemize}

\begin{tcolorbox}[breakable,colback=teal!4,colframe=teal!35!black,title={Before you move on}]
What does die Komödie mean? (\textquotedblleft{}Comedy\textquotedblright{}.) Say it once more.
\end{tcolorbox}
\section[die Tragödie]{die Tragödie — tragedy}

\label{lesson:GE-C524-tragoedie}



\begin{quote}
Before the new one: say the German for ratings, then the German for comedy.
\end{quote}

\subsection*{You'll want to know: die Tragödie}
\textbf{die Tragödie} — \textquotedblleft{}tragedy\textquotedblright{}.

\begin{quote}
Das Theater spielt eine Tragödie aus der Antike.
\end{quote}

\begin{quote}
\emph{The theatre is staging a tragedy from ancient times.}
\end{quote}

Plural: \emph{die Tragödien}. Also a real disaster.

\begin{cousinweb}
\emph{Tragödie} comes from Latin \emph{tragoedia}, from Greek \emph{tragōidía}, built on \emph{trágos} \textquotedblleft{}male goat\textquotedblright{} and \emph{ōid\'{\={e}}} \textquotedblleft{}song\textquotedblright{} (EtymDB).
\end{cousinweb}

\begin{grammarlens}[title={What you've built}]
One more word for telling, explaining and describing at B1.
\end{grammarlens}

\subsection*{Your turn}
\begin{itemize}
\raggedright
  \item \emph{Say it:} \emph{die Tragödie}
  \item \emph{Say it:} \emph{die Tragödie}, once more
  \item \emph{Say it:} a sentence of your own with \emph{die Tragödie}
  \item \emph{Recall:} say the German for ratings, then the German for comedy, then say \emph{die Tragödie} again
\end{itemize}

\begin{tcolorbox}[breakable,colback=teal!4,colframe=teal!35!black,title={Before you move on}]
What does die Tragödie mean? (\textquotedblleft{}Tragedy\textquotedblright{}.) Say it once more.
\end{tcolorbox}
\section[der Dokumentarfilm]{der Dokumentarfilm — documentary}

\label{lesson:GE-C524-dokumentarfilm}



\begin{quote}
Before the new one: say the German for comedy, then the German for tragedy.
\end{quote}

\subsection*{You'll want to know: der Dokumentarfilm}
\textbf{der Dokumentarfilm} — \textquotedblleft{}documentary\textquotedblright{}.

\begin{quote}
Im Fernsehen läuft ein Dokumentarfilm über Wale.
\end{quote}

\begin{quote}
\emph{There's a documentary about whales on TV.}
\end{quote}

Plural: \emph{die Dokumentarfilme}. In speech also \emph{die Doku}.

\begin{cousinweb}
\emph{Dokumentarfilm} is built on \emph{Dokument}, from Latin \emph{documentum} \textquotedblleft{}proof, lesson\textquotedblright{}, and \emph{Film}, borrowed from English \emph{film} (EtymDB).
\end{cousinweb}

\begin{grammarlens}[title={What you've built}]
One more word for telling, explaining and describing at B1.
\end{grammarlens}

\subsection*{Your turn}
\begin{itemize}
\raggedright
  \item \emph{Say it:} \emph{der Dokumentarfilm}
  \item \emph{Say it:} \emph{der Dokumentarfilm}, once more
  \item \emph{Say it:} a sentence of your own with \emph{der Dokumentarfilm}
  \item \emph{Recall:} say the German for comedy, then the German for tragedy, then say \emph{der Dokumentarfilm} again
\end{itemize}

\begin{tcolorbox}[breakable,colback=teal!4,colframe=teal!35!black,title={Before you move on}]
What does der Dokumentarfilm mean? (\textquotedblleft{}Documentary\textquotedblright{}.) Say it once more.
\end{tcolorbox}
\section[der Untertitel]{der Untertitel — subtitle}

\label{lesson:GE-C524-untertitel}



\begin{quote}
Before the new one: say the German for tragedy, then the German for documentary.
\end{quote}

\subsection*{You'll want to know: der Untertitel}
\textbf{der Untertitel} — \textquotedblleft{}subtitle\textquotedblright{}.

\begin{quote}
Ich schaue Filme lieber mit Untertiteln.
\end{quote}

\begin{quote}
\emph{I prefer watching films with subtitles.}
\end{quote}

Plural: \emph{die Untertitel}.

\begin{cousinweb}
\emph{Untertitel} is built on \emph{unter} \textquotedblleft{}under\textquotedblright{} and \emph{Titel} \textquotedblleft{}title\textquotedblright{}, which goes back to Latin \emph{titulus} \textquotedblleft{}inscription, label\textquotedblright{} (Lewis \& Short).
\end{cousinweb}

\begin{grammarlens}[title={What you've built}]
One more word for telling, explaining and describing at B1.
\end{grammarlens}

\subsection*{Your turn}
\begin{itemize}
\raggedright
  \item \emph{Say it:} \emph{der Untertitel}
  \item \emph{Say it:} \emph{der Untertitel}, once more
  \item \emph{Say it:} a sentence of your own with \emph{der Untertitel}
  \item \emph{Recall:} say the German for tragedy, then the German for documentary, then say \emph{der Untertitel} again
\end{itemize}

\begin{tcolorbox}[breakable,colback=teal!4,colframe=teal!35!black,title={Before you move on}]
What does der Untertitel mean? (\textquotedblleft{}Subtitle\textquotedblright{}.) Say it once more.
\end{tcolorbox}
\section[synchronisieren]{synchronisieren — to dub; to synchronise}

\label{lesson:GE-C524-synchronisieren}



\begin{quote}
Before the new one: say the German for documentary, then the German for subtitle.
\end{quote}

\subsection*{You'll want to know: synchronisieren}
\textbf{synchronisieren} — \textquotedblleft{}to dub; to synchronise\textquotedblright{}.

\begin{quote}
Bei uns werden fast alle Filme synchronisiert.
\end{quote}

\begin{quote}
\emph{Here almost all films are dubbed.}
\end{quote}

Regular, Perfekt without \emph{ge-}. Also to sync two devices.

\begin{cousinweb}
\emph{synchronisieren} is built on \emph{synchron} \textquotedblleft{}at the same time\textquotedblright{}, a learned word from Greek \emph{sýnchronos}, built on \emph{khrónos} \textquotedblleft{}time\textquotedblright{}; EtymDB has no entry for it.
\end{cousinweb}

\begin{grammarlens}[title={What you've built}]
One more word for telling, explaining and describing at B1.
\end{grammarlens}

\subsection*{Your turn}
\begin{itemize}
\raggedright
  \item \emph{Say it:} \emph{synchronisieren}
  \item \emph{Say it:} \emph{synchronisieren}, once more
  \item \emph{Say it:} a sentence of your own with \emph{synchronisieren}
  \item \emph{Recall:} say the German for documentary, then the German for subtitle, then say \emph{synchronisieren} again
\end{itemize}

\begin{tcolorbox}[breakable,colback=teal!4,colframe=teal!35!black,title={Before you move on}]
What does synchronisieren mean? (\textquotedblleft{}To dub; to synchronise\textquotedblright{}.) Say it once more.
\end{tcolorbox}
